% Abstract — ~250 words.
Diabetic retinopathy (DR) is a leading cause of preventable blindness, and
automated screening systems are expected to scale early detection. While
recent deep-learning models for DR classification report high accuracies,
clinical deployment is hindered by their tendency to produce
\emph{over-confident} predictions: when a model is wrong, it is often
just as confident as when it is right.

This thesis builds an \emph{uncertainty-aware} DR grading system on top of
the APTOS~2019 fundus dataset. Six baseline architectures (ResNet50,
DenseNet121, VGG16, Xception, a custom CNN, and a CNN with mixed
\textit{tanh}+\textit{ReLU} activations) are first evaluated as binary
classifiers using a stratified 70/15/15 split. With proper preprocessing
and class weighting, all five strong baselines reach 95--96\,\%
test accuracy, which McNemar tests show is statistically indistinguishable
across them. A heterogeneous ensemble pushes accuracy to 96.55\,\% and
predictive entropy enables \emph{selective prediction} that reaches
99.6\,\% accuracy at 50\,\% coverage.

We then add three uncertainty-aware components: (i) split conformal
prediction with formal coverage guarantees at $\alpha = 0.10$ and
$\alpha = 0.05$; (ii) Monte Carlo Dropout retrained variants of the CNN
and ResNet50; and (iii) multi-method out-of-distribution detection in
which feature-space methods (Mahalanobis, cosine) achieve perfect
separation against synthetic OOD inputs. Five-fold cross-validation
on ResNet50 gives a tight test accuracy of $95.64 \pm 0.18$\,\%, ruling out
single-split luck.

Finally, we re-formulate the task as the clinically meaningful 5-stage
grading scale (No DR, Mild, Moderate, Severe, PDR). The ResNet50
multi-class model achieves a quadratic-weighted kappa of 0.847,
competitive with public APTOS~2019 leaderboards. Combined with conformal
prediction at 90\,\% coverage, the system produces interpretable
prediction sets and a principled clinical refer-to-specialist rule.
